% TODO make hw stylesheet
\documentclass[10pt]{scrartcl}
\usepackage[a4paper, total={6in, 8in}]{geometry}
\usepackage[usenames,dvipsnames,svgnames]{xcolor}
\usepackage[shortlabels]{enumitem}
\usepackage[framemethod=TikZ]{mdframed}
\usepackage{amsmath,amssymb,amsthm}
\usepackage[colorlinks]{hyperref}
\usepackage{multicol}
\usepackage{tabularx}

\hypersetup{
	colorlinks=true,
	linkcolor=blue,
	filecolor=magenta,      
	urlcolor=magenta,
	pdftitle={MAT 322 HW}
}

\usepackage{microtype}
\usepackage{mathtools}
\usepackage[headsepline]{scrlayer-scrpage}
\usepackage{thmtools}
\usepackage[textsize=footnotesize]{todonotes}

\mdfdefinestyle{mdbluebox}{roundcorner=10pt,innerbottommargin=9pt,
	linecolor=blue,backgroundcolor=TealBlue!5,}
\declaretheoremstyle[headfont=\sffamily\bfseries\color{MidnightBlue},
mdframed={style=mdbluebox},]{thmbluebox}

\mdfdefinestyle{mdbloobox}{frametitlefont=\bfseries,innerbottommargin=8pt,
	nobreak=true,backgroundcolor=TealBlue!5,linecolor=blue,}
\declaretheoremstyle[headfont=\bfseries\color{MidnightBlue},
mdframed={style=mdbloobox},headpunct={\\[3pt]},postheadspace=0pt,]{thmbloobox}

\mdfdefinestyle{mdredbox}{frametitlefont=\bfseries,innerbottommargin=8pt,
	nobreak=true,backgroundcolor=Salmon!5,linecolor=RawSienna,}
\declaretheoremstyle[headfont=\bfseries\color{RawSienna},
mdframed={style=mdredbox},headpunct={\\[3pt]},postheadspace=0pt,]{thmredbox}

\mdfdefinestyle{mdgreenbox}{linecolor=ForestGreen,backgroundcolor=ForestGreen!5,
	linewidth=2pt,rightline=false,leftline=true,topline=false,bottomline=false,}
\declaretheoremstyle[headfont=\bfseries\sffamily\color{ForestGreen!70!black},
mdframed={style=mdgreenbox},headpunct={ --- },]{thmgreenbox}

\mdfdefinestyle{mdorangebox}{linecolor=RedOrange,backgroundcolor=RedOrange!5,
	linewidth=2pt,rightline=false,leftline=true,topline=false,bottomline=false,}
\declaretheoremstyle[headfont=\bfseries\sffamily\color{RedOrange!70!black},
mdframed={style=mdorangebox},headpunct={ --- },]{thmorangebox}

\mdfdefinestyle{mdblackbox}{linecolor=black,backgroundcolor=RedViolet!5!gray!5,
	linewidth=3pt,rightline=false,leftline=true,topline=false,bottomline=false,}
\declaretheoremstyle[mdframed={style=mdblackbox}]{thmblackbox}

\declaretheoremstyle[spaceabove=6pt,spacebelow=6pt]{basehead}

\declaretheorem[style=basehead,name=Problem]{problem}
\declaretheorem[style=basehead,name=Problem,numbered=no]{problem*}
\declaretheorem[style=thmbloobox,name=Setup]{setup} % for config geo
\declaretheorem[style=thmredbox,name=Example,sibling=problem]{example}
\declaretheorem[style=thmredbox,name=$\bigstar$ Problem,sibling=problem]{niceprob}
\declaretheorem[style=thmbluebox,name=Theorem,sibling=problem]{theorem}
\declaretheorem[style=thmbluebox,name=Corollary,sibling=theorem]{corollary}
\declaretheorem[style=thmbluebox,name=Proposition,sibling=theorem]{prop}
\declaretheorem[style=thmbluebox,name=Lemma,numberwithin=problem]{lemma}
\declaretheorem[style=thmbluebox,name=Theorem,numbered=no]{theorem*}
\declaretheorem[style=thmbluebox,name=Lemma,numbered=no]{lemma*}
\declaretheorem[style=thmgreenbox,name=Claim,numberwithin=problem]{claim}
\declaretheorem[style=thmgreenbox,name=Claim,numbered=no]{claim*}
\declaretheorem[style=thmblackbox,name=Remark,numberwithin=problem]{remark}
\declaretheorem[style=thmblackbox,name=Remark,numbered=no]{remark*}
\declaretheorem[style=thmgreenbox,name=Definition,sibling=theorem]{definition}
\declaretheorem[style=thmgreenbox,name=Definition,numbered=no]{definition*}
\declaretheorem[style=thmorangebox,name=Warning,sibling=theorem]{warning}
\declaretheorem[style=thmorangebox,name=Warning,numbered=no]{warning*}
\declaretheorem[style=thmorangebox,name=Note,numbered=no]{note}
\declaretheorem[style=thmblackbox,name=Exercise,sibling=problem]{exercise}
\declaretheorem[style=thmblackbox,name=Exercise,numbered=no]{exercise*}
\declaretheorem[style=thmblackbox,name=Question,sibling=problem]{question}
\declaretheorem[style=thmblackbox,name=Question,numbered=no]{question*}

\addtolength{\textheight}{3.14cm}
\providecommand{\ol}{\overline}
\providecommand{\ul}{\underline}
\providecommand{\eps}{\varepsilon}
\providecommand{\half}{\frac{1}{2}}
\providecommand{\dang}{\measuredangle} %% Directed angle
\providecommand{\CC}{\mathbb C}
\providecommand{\FF}{\mathbb F}
\providecommand{\NN}{\mathbb N}
\providecommand{\QQ}{\mathbb Q}
\providecommand{\RR}{\mathbb R}
\providecommand{\ZZ}{\mathbb Z}
\providecommand{\dg}{^\circ}
\providecommand{\ii}{\item}
\providecommand{\setdiff}{\smallsetminus}
\providecommand{\alert}[1]{{\sffamily\textbf{\textcolor{blue}{#1}}}}
\providecommand{\inv}{^{-1}}
\providecommand{\mb}[1]{\mathbf{#1}}
\providecommand{\dif}{\;d}

\reversemarginpar

\newenvironment{soln}{\begin{proof}[Solution]}{\end{proof}}
\newenvironment{walkthrough}{\noindent \textbf{Walkthrough.}}{\medskip}
\newlist{walk}{enumerate}{3}
\setlist[walk]{label=\bfseries (\alph*)}

\providecommand{\pow}{\operatorname{Pow}}
\providecommand{\vocab}[1]{{\textbf{\textcolor{ForestGreen}{#1}}}}
\providecommand{\lcm}{\operatorname{lcm}}
\providecommand{\abs}[1]{\left|#1\right|}

\usepackage{asymptote}
\begin{asydef}
	defaultpen(fontsize(10pt));
	unitsize(1cm);
	size(8cm); // set a reasonable default
	usepackage("amsmath");
	usepackage("amssymb");
	settings.tex="pdflatex";
	settings.outformat="pdf";
	// Replacement for olympiad+cse5 which is not standard
	import geometry;
	// recalibrate fill and filldraw for conics
	void filldraw(picture pic = currentpicture, conic g, pen fillpen=defaultpen, pen drawpen=defaultpen)
	{ filldraw(pic, (path) g, fillpen, drawpen); }
	void fill(picture pic = currentpicture, conic g, pen p=defaultpen)
	{ filldraw(pic, (path) g, p); }
	// some geometry
	pair foot(pair P, pair A, pair B) { return foot(triangle(A,B,P).VC); }
	pair orthocenter(pair A, pair B, pair C) { return orthocentercenter(A,B,C); }
	pair centroid(pair A, pair B, pair C) { return (A+B+C)/3; }
	// cse5 abbreviations
	path CP(pair P, pair A) { return circle(P, abs(A-P)); }
	path CR(pair P, real r) { return circle(P, r); }
	pair IP(path p, path q) { return intersectionpoints(p,q)[0]; }
	pair OP(path p, path q) { return intersectionpoints(p,q)[1]; }
	path Line(pair A, pair B, real a=0.6, real b=a) { return (a*(A-B)+A)--(b*(B-A)+B); }
	// cse5 more useful functions
	picture CC() {
		picture p=rotate(0)*currentpicture;
		currentpicture.erase();
		return p;
	}
	pair MP(Label s, pair A, pair B = plain.S, pen p = defaultpen) {
		Label L = s;
		L.s = "$"+s.s+"$";
		label(L, A, B, p);
		return A;
	}
	pair Drawing(Label s = "", pair A, pair B = plain.S, pen p = defaultpen) {
		dot(MP(s, A, B, p), p);
		return A;
	}
	path Drawing(path g, pen p = defaultpen, arrowbar ar = None) {
		draw(g, p, ar);
		return g;
	}
\end{asydef}

\usepackage{mathrsfs}

\ihead{\footnotesize MAT 314 HW07}
\ohead{\footnotesize Last compiled \today}

\title{\vspace{-2em}MAT 314 HW07}
\author{\large Aditya Pahuja}
\date{\large Due 04/29}
\begin{document}
\maketitle

\begin{problem*}1.1.
	Let $R$ be an integral domain that contains a field $F$
	as a subring and that is finite-dimensional when viewed as
	a vector space over $F$. Prove that $R$ is a field.
\end{problem*}

\begin{soln}
	We show that any nonzero $r\in R$ has an inverse.
	If $n$ is the dimension of the $F$-vector space $R$, then
	$1$, $r$, \dots, $r^n$ is linearly dependent, satisfying some relation
	$c_0 + c_1r + \cdots + c_nr^n = 0$. Let $k\ge 0$ be the smallest integer
	for which $c_k\ne 0$; obviously $k<n$. Since $R$ is an integral domain
	we can deduce that
	\[c_k + c_{k+1}r + \cdots + c_nr^{n-k} = 0
	\iff r\cdot (-c_k\inv)(c_{k+1} + c_{k+2}r + \cdots + c_n r^{n-k-1}) = 1\]
	so that $r$ has an inverse (this exact manipulation is also used
	in the following problem).
\end{soln}

\begin{problem*}2.2.
	Let $f(x) = x^n - a_{n-1}x^{n-1} + \cdots + (-1)^n a_0$ be
	an irreducible polynomial over $F$, and let $\alpha$ be a root of $f$
	in the extension field $K$. Determine the element $\alpha\inv$ explicitly
	in terms of $\alpha$ and the coefficients $a_i$.
\end{problem*}

\begin{soln}
	Since $f$ is irreducible over $F$, $a_0$ must be nonzero.
	Then, $f(\alpha)=0$ so
	\[-\frac{(-1)^n}{a_0}(f(\alpha) - (-1)^n a_0) = 1.\]
	so the inverse of $\alpha$ is $-(-1)^na_0\inv\cdot
	(\alpha^{n-1} - a_{n-1}\alpha^{n-2} + \cdots + (-1)^{n-1}a_1)$,
	factoring out $\alpha$ from $f(\alpha) - (-1)^na_0$.
\end{soln}

\begin{problem*}3.1.
	Let $F$ be a field, and let $\alpha$ be an element that generates
	a field extension of $F$ of degree $5$.
	Prove that $\alpha^2$ generates the same extension.
\end{problem*}

\begin{soln}
	Obviously $\alpha^2\in F(\alpha)$, so $F(\alpha^2)\subseteq F(\alpha)$.
	On the other hand,
	\[[F(\alpha):F(\alpha^2)][F(\alpha^2):F] = 5\]
	and $\alpha$ is a root of the quadratic polynomial
	$p(X) = X^2 - \alpha^2\in F(\alpha^2)[X]$, so the degree of $F(\alpha)$
	over $F(\alpha^2)$ is either $1$ or $2$; being a divisor of $5$,
	the degree must be $1$, and so $F(\alpha) = F(\alpha^2)$.
\end{soln}

\begin{problem*}8.1.
	Prove that every finite extension of a finite field has a primitive element.
\end{problem*}

\begin{soln}
	Let $K$ be a finite extension of the finite field $\FF_q$,
	with $q = p^m$ and $|K| = p^n$ for prime $p$ and positive integers $m$, $n$.
	Then, if $\alpha$ generates the cyclic group $K^\times$,
	we can see that $\FF_q(\alpha)$ on one hand contains $K$
	(since it contains $K^\times$, of course)
	but on the other hand is contained in $K$ because $\FF_q\subseteq K$
	and $\alpha\in K$, so $\alpha$ is a primitive element for the extension.
\end{soln}

\begin{problem*}10.1.
	Prove that the subset of $\CC$ consisting of the algebraic numbers
	is algebraically closed.
\end{problem*}

\begin{soln}
	Let $\ol\QQ$ be the set of algebraic numbers
	and let $p(x) = x^n + a_{n-1}x^{n-1} + \cdots + a_0$
	be a polynomial with coefficients in $\ol\QQ$. Then,
	$p$ has coefficients in the finite extension $K=\QQ(a_0,\dots,a_{n-1})$
	of $\QQ$. We wish to show that the roots of $p$ are also algebraic numbers,
	i.e. are contained in a finite extension of $\QQ$.
	If $\beta$ is any root of $p$, then $K(\beta)= \QQ(a_0,\dots,a_{n-1},\beta)$
	is spanned by $\{v_i\beta^j : 0\le i< m,\;0\le\beta < n\}$
	over $\QQ$ where $\{v_1, \dots, v_m\}$ is a basis for $K$.
	Thus $K(\beta)$ is a finite extension of $\QQ$,
	so $\beta$ is an algebraic number, hence any root of a polynomial
	with algebraic number coefficients is also algebraic.
\end{soln}




















\end{document}